\section{Conclusion}
\label{sec:conclusion}

% (~1-2 pages)

% Summary of key findings:
%   1. Persistent homology reveals distinct topological signatures in
%      microbial co-occurrence networks associated with neurotransmitter
%      production capacity (serotonin, GABA, dopamine, SCFAs)
%   2. NT-producing community topology — particularly persistent H$_1$ loops
%      representing cross-feeding cycles — predicts NT production potential
%      better than taxonomic composition alone
%   3. Mycobiome (fungal) co-colonization fragments NT-network topology,
%      especially the GABA and serotonin pathways, providing a topological
%      mechanism for mold-associated neurological symptoms
%   4. AMR carriers act as topological disruptors with a quantifiable
%      ``tipping point'' beyond which NT cross-feeding networks collapse
%   5. Topological regime transitions precede compositional shifts in
%      NT-associated taxa, supporting topology as an early warning system

% Contributions:
%   - First application of persistent homology (not Mapper) to microbiome
%     co-occurrence networks with Betti curves and persistence landscapes
%   - First linking of topological features to a functional outcome
%     (neurotransmitter production)
%   - First topological characterization of mycobiome-bacterial and
%     AMR-bacterial disruption
%   - Bridges TDA-microbiome and TDA-neuroscience literatures

% Broader implications:
%   - Gut-brain axis: community topology, not just composition,
%     determines NT signaling capacity
%   - Clinical: topological monitoring during antibiotic treatment,
%     probiotic interventions, environmental exposure management
%   - Personalized medicine: topological profiles could guide
%     targeted interventions to restore NT-producing network structure

% Closing:
%   The gut microbiome is not merely a collection of NT-producing species ---
%   it is a structured ecological network whose topology determines its
%   collective capacity for neurotransmitter biosynthesis. By revealing
%   how this topology is disrupted by fungi, antibiotics, and resistant
%   organisms, and how it correlates with neurotransmitter output and
%   neurological outcomes, we open a new structural avenue for understanding
%   and optimizing the gut-brain axis.
